We now specialise to the order flow of Chapter \ref{chap.orderDrivenMarkets}.
The six event types are the market, limit and withdrawal orders on each side:
\begin{center}
 \begin{tabular}{lcccccc}
  \toprule
  $e$ & market buy & market sell & limit buy & limit sell & withdraw bid & withdraw ask \\
  \midrule
  $\pressure_e$ & $+1$ & $-1$ & $+1$ & $-1$ & $-1$ & $+1$ \\
  \bottomrule
 \end{tabular}
\end{center}
The vector $\pressure$ records the direction in which each type pushes the mid price:
a market buy consumes the ask and a withdrawal from the bid removes support,
so the first pushes up and the second down.

\begin{defi}\label{def.intensityContrast}
 Let $\pressure \in \lbrace -1, 0, +1\rbrace^{\numEventTypes}$, and set
 \begin{equation*}
  \upIntensity(t) := \sum_{e \, : \, \pressure_e = +1} \intensity(t),
  \qquad
  \downIntensity(t) := \sum_{e \, : \, \pressure_e = -1} \intensity(t) .
 \end{equation*}
 The intensity contrast is
 \begin{equation}\label{eq.intensityContrast}
  \intensityContrast(t) := \upIntensity(t) - \downIntensity(t)
  = \pressure^{\transpose}\intensity[ ](t)
  = \pressure^{\transpose}\baseIntensity
  + \pressure^{\transpose}\excitation \decayedCounts(t) .
 \end{equation}
\end{defi}

The restriction of $\pressure$ to $\lbrace -1,0,+1\rbrace$ is what makes the middle equality
hold: for a general real vector $\pressure^{\transpose}\intensity[ ]$ is a weighted contrast
and is not the difference of two sums, and $\upIntensity$, $\downIntensity$ have no meaning.

The last term of \eqref{eq.intensityContrast} is the one that carries information,
and the first is a constant.
Under the symmetry
\begin{equation}\label{eq.baselineSymmetry}
 \pressure^{\transpose}\baseIntensity = 0,
\end{equation}
that is, when the up-pushing and down-pushing baselines balance,
the contrast reduces to $\pressure^{\transpose}\excitation\decayedCounts$ exactly,
and its sign is a statement about the state alone.
Every parametrisation used in these notes satisfies \eqref{eq.baselineSymmetry},
by construction: the baselines are equal in pairs.
We state the assumption here, where it is used, rather than folding it into the definition,
because it is not innocuous --- on an asymmetric parametrisation the contrast has a constant
offset, and a rule that trades on its sign inherits that offset as a permanent bias.

Note also that the $\eone$-th entry of $\pressure^{\transpose}\excitation$
is the signed jump in $\intensityContrast$ caused by one event of type $\eone$,
and that it scales with $\decay$ at fixed $\branchingRatio$,
whereas $\pressure^{\transpose}\branchingMatrix$ does not.

\subsection{The contrast and the order flow imbalance}

The contrast $\intensityContrast$ is an oracle: it is a function of the model's parameters and
of its state, both of which we know only because we generated the path.
The order flow imbalance of Chapter \ref{chap.orderDrivenMarkets},
$\OFI_{t,w} = \sum_n \orderFlowContribution$ over the events in $(t-w, t]$,
is an estimator built from observables alone.
It is tempting to describe the second as a crude version of the first --- a causal linear
filter on the same state, with a boxcar in place of the exponential.
The description is wrong, and it is worth saying why, because it is the difference that makes
$\OFI$ worth computing.

The two are statistics of different things.

$\intensityContrast$ sees the self-excited part of the intensity, and only that.
By \eqref{eq.intensityContrast} it is a function of $\decayedCounts$,
so it knows nothing about the immigrants that are about to arrive;
under \eqref{eq.baselineSymmetry} it is blind to $\baseIntensity$ altogether.
It also requires $\excitation$ and $\decay$ to be known.

$\OFI$ is model free.
It counts events as they occurred, immigrants and offspring alike,
without asking which was which,
and it therefore carries information about the exogenous flow that $\intensityContrast$
discards by construction.

And $\OFI$ reads the book.
The contribution $\orderFlowContribution$ of a single event is a touch-size increment:
it carries the size of the order and the state of the queue it arrived at,
neither of which appears anywhere in $\hawkesKernel$.
The kernel governs \emph{when} events happen and of \emph{which type};
the marks govern \emph{how much} they move.
No function of $\decayedCounts$ can recover a mark.

There remain three differences of shape, and they are real, but they are differences between
two estimators of different quantities rather than defects of one of them:
the window is a boxcar where the kernel is an exponential;
the types are weighted by $\pressure$ alone where the kernel weights them by
$\pressure^{\transpose}\excitation$;
and the marks enter one and not the other.

\subsection{What the model does not contain}

The apparatus of this chapter generates order flow in a controlled setting,
and a reader should not mistake a result obtained in it for a result about a traded market.
The following are absent by construction.

The intensities do not read the book: there is no queue-depletion feedback,
so the arrival rates are indifferent to whether a queue is about to empty.
The \emph{marks} do read it, and that is the mechanism behind the book's restoring force.
There is no informed trading and no adverse selection,
which is where the value of the queue imbalance $\queueImbalance[1]$ would otherwise live.
There is no queue-position economics.
There is a single excitation timescale, the kernel being one exponential,
per Remark \ref{remark.stateSize}.
There is no intraday non-stationarity, so a window tuned here is not tested against one.
And there is no anchor for the price level, so the level wanders freely
and the mean price change over a session is a per-seed random drift.

Two questions cannot be asked in this setting at all,
because they require recorded data:
whether any of the regimes studied here occurs in a traded market,
and whether one excitation timescale suffices to describe one.
